\section{Background and Problem Setting}
\label{sec:background}

%TL;DR of paragraph: notation for a structure and an arm, and the two assumptions that separate this setting from image generation.
A \emph{structure} here is a map $x : \{1,\dots,d\}^3 \to \mathcal{V}$ from a voxel grid to a
categorical block vocabulary $\mathcal{V}$, with one symbol reserved for air. An
\emph{arm} is a finite set of structures with a name and a provenance record: a
generative model's samples, a procedural generator's output, or real data under a known
corruption. Evaluation compares an arm's empirical distribution $Q$ against a fixed
reference $P$ of real structures. Two features of this setting differ from image
generation. Sample sizes are in the hundreds rather than the tens of thousands, and arms
may differ in size, so a statistic whose null depends on $n$ ranks by sample size.

%TL;DR of paragraph: define MMD, the unbiased estimator, and how KID and FID instantiate it on features.
Both metrics we report are kernel two-sample distances. For a feature map $\phi$ and a
kernel $k$,
\begin{equation}
\mathrm{MMD}^2(P,Q) = \mathbb{E}\,k(\phi(x),\phi(x')) + \mathbb{E}\,k(\phi(y),\phi(y'))
  - 2\,\mathbb{E}\,k(\phi(x),\phi(y)),
\label{eq:mmd}
\end{equation}
with $x,x' \sim P$ and $y,y' \sim Q$ \citep{gretton2012mmd}. The unbiased estimator of
Eq.~\ref{eq:mmd} omits the diagonal terms $k(\phi(x_i),\phi(x_i))$ from the within-set
means. The biased estimator keeps them, so its null on two disjoint real halves decays
with $n$. KID \citep{binkowski2018kid} is the unbiased estimator with a polynomial
kernel on encoder features, and FID \citep{heusel2017fid} fits a Gaussian to each feature
set instead. A \emph{render-and-embed} metric is any of these with $\phi =
(\text{encoder}) \circ (\text{render})$, so the structure reaches the statistic only as a
set of images. 

%TL;DR of paragraph: define the noise floor, the unit that makes metrics on different features commensurable.
Raw distances across features can differ by orders of magnitude, so we report all distances in \emph{noise floors}. For each feature and sample size $n$, the noise floor is defined as the standard deviation of the distance obtained from repeated real-vs-reference draws under the same experimental run. Thus, a value of $1.0$ indicates that the observed distance is one standard deviation of the metric's sampling variability. Because the sampling variability changes with $n$, the noise-floor unit is $n$-dependent (Appendix~\ref{app:nscale}). \com{check all appendix refs}

%TL;DR of paragraph: define the three validity properties a metric must have, which we later apply as blocking gates.
We treat a metric as \emph{valid} for ranking when it has three properties.
\begin{enumerate}
    \item 
\textbf{Ordering}: its value rises monotonically with the severity of a known corruption.
% its value rises monotonically with the severity of a corruption known to reduce quality. 
    \item \textbf{Invariance}: its value does not move under a transformation
that leaves quality unchanged. 
\item \textbf{$n$-stability}: its null on real-vs-real draws
does not drift with $n$. 
\end{enumerate}
Sensitivity, how \emph{much} damage a metric
resolves at a given $n$, is statistical power rather than measuring the right quantity, so
we record it alongside the three rather than as a fourth.
